\chapter{Топологический дефицит замыкания и неизбежность дефекта}

Исходная онтология проекта утверждала, что материя не добавлена к вакууму,
а возникает как остаток невозможности его идеального самозамыкания. Ранний
язык ``неплотной упаковки'' был эвристическим. Теперь уже построенный
оператор Тёплица позволяет дать этому тезису строгую операторную форму,
не обращаясь к упаковкам сфер или решётке Лича.

\section{Совершенное замыкание как обратимость}

Пусть $q_0\in M_{105}(\mathbb C)$ --- канонический проектор ранга
пятнадцать, а $S$ --- односторонний сдвиг на $H^2(S^1)$. На полном
коэффициентном пространстве определим
\begin{equation}
 \widehat T_+
 =S\otimes q_0+I\otimes(1-q_0).
 \label{eq:v5-full-coefficient-toeplitz}
\end{equation}
Тогда
\begin{equation}
 \widehat T_+^*\widehat T_+=I,
 \qquad
 I-\widehat T_+\widehat T_+^*
 =|e_0\rangle\langle e_0|\otimes q_0.
 \label{eq:v5-one-sided-closure-defect}
\end{equation}
Следовательно, оператор является изометрией, но не унитарен: после одного
полного шага остаётся коядро размерности пятнадцать.

Совершенным замыканием будем называть обратимый фредгольмов оператор. Для
него одновременно
\begin{equation}
 \ker F=0,
 \qquad
 \operatorname{coker}F=0.
\end{equation}
Это точный аналог состояния, которое после композиции переходов полностью
возвращается в себя и не оставляет локализованного остатка.

\section{Нижняя граница неплотности}

Для фредгольмова оператора $F$ положим
\begin{equation}
 \Delta_{\rm cl}(F)
 =\frac{\dim\ker F+\dim\operatorname{coker}F}{105}.
 \label{eq:v5-closure-deficit-definition}
\end{equation}
Если $F$ принадлежит ориентированному классу $\widehat T_+$, то
\begin{equation}
 \operatorname{ind}F
 =\dim\ker F-\dim\operatorname{coker}F=-15.
\end{equation}
Для любых неотрицательных целых $k,c$ выполнено
\begin{equation}
 k+c\geq |k-c|.
\end{equation}
Поэтому
\begin{equation}
 \Delta_{\rm cl}(F)\geq\frac{15}{105}=\frac17.
 \label{eq:v5-closure-deficit-bound}
\end{equation}
Равенство достигается тогда и только тогда, когда отсутствуют добавочные
пары ядро--коядро. Явный оператор \eqref{eq:v5-full-coefficient-toeplitz}
имеет $k=0$, $c=15$ и достигает этой границы.

\begin{theorem}[Минимальный неизбежный дефект]
В топологическом секторе $\operatorname{ind}F=-15$ совершенное замыкание
невозможно. Любой представитель имеет не менее пятнадцати незамкнутых
каналов, а канонический оператор Тёплица реализует ровно пятнадцать.
Нормированный минимальный дефицит равен $1/7$.
\end{theorem}

Индекс устойчив при компактных возмущениях. Поэтому локальная перестройка
оператора может перемещать дефект и создавать дополнительные пары
нулевых мод, но не может удалить минимальный остаток без выхода из
топологического класса. В языке соответствия объём--граница это стандартная
роль граничной карты расширения; см. \path{arXiv:1411.7527} и
\path{arXiv:1604.02337}.

\section{Почему Real-пара не уничтожает дефект}

Сопряжённая ориентация имеет индекс $+15$. Обычные комплексные индексы
пары складываются в нуль, но это не означает совершенного замыкания:
каждая половина имеет дефектный проектор ранга пятнадцать. Для полного
обменного пакета естественная нормировка даёт
\begin{equation}
 \Delta_{\rm cl}^{\rm Real}
 =\frac{15+15}{105+105}=\frac17.
 \label{eq:v5-real-pair-closure-deficit}
\end{equation}

Более существенно, обменная пара представляет ненулевой класс
\begin{equation}
 15\in KO_6\bigl(M_{105}(\mathbb C)_{\mathbb R}\bigr).
\end{equation}
Симметрийно-сохраняющая
гомотопия к обратимому нулевому представителю обнулила бы этот класс, что
противоречит уже доказанной инъективной комплексификации
$c_6(15)=(-15,+15)$. Необходимость вещественных алгебр и $KKO$ при
антилинейных симметриях является стандартной частью bulk--edge
соответствия; см. \path{arXiv:1604.02337}. Явные Real-граничные карты
даны в \path{arXiv:1504.03284}.

\begin{remark}[Что именно сокращается]
Сокращается ориентированный знак обычного комплексного индекса. Не
сокращаются два положительных дефектных проектора и обменный Real-класс.
Поэтому нулевой суммарный комплексный индекс нельзя интерпретировать как
отсутствие материи.
\end{remark}

\section{Совпадение топологического и вариационного минимумов}

Ранее независимо было доказано, что минимальная энергия внутренней петли
равна
\begin{equation}
 S_{\rm loop}(V_{15})=\frac{30}{210}=\frac17.
\end{equation}
Теперь видно, почему это число повторяется. Целочисленная выпуклость
распределяет winding по пятнадцати единичным каналам, а индексная нижняя
граница требует не менее пятнадцати дефектных каналов. Канонический
$V_{15}$ одновременно минимизирует изменение петли и насыщает нижнюю
границу незамыкания.

Это первое строгое совпадение ранней онтологии с поздней операторной
конструкцией:
\begin{equation}
 \begin{aligned}
  \text{минимальная стоимость внутреннего перехода}
  &=\text{минимальная доля незамыкания}\\
  &=\frac17.
 \end{aligned}
 \label{eq:v5-two-minima-one-seventh}
\end{equation}
Равенство пока безразмерно. Оно не превращает $1/7$ в массу и не задаёт
пространственный радиус.

\section{Граница доказанного постулата}

Полученный результат имеет условие: он доказывает неизбежность дефекта
\emph{в уже выбранном ненулевом классе пятнадцать}. Он ещё не доказывает,
что вакуум обязан выбирать именно этот сектор вместо нулевого класса, и
не классифицирует все частицы как его дефекты.

Кроме того, $\Delta_{\rm cl}$ является индексной плотностью, а не
геометрической долей незаполненного объёма. Буквальная подстановка
коэффициентов плотнейшей упаковки запрещена: строгая ветвь проекта уже
отделила число $24$ от ранних решёточных метафор.

\begin{nogocriterion}[Запрет объявления полной теории материи]
Из нижней границы $\Delta_{\rm cl}\geq1/7$ нельзя заключать, что выведены
масса, размер или спектр всех частиц. Для этого один родительский
функционал должен сначала выбрать ненулевой сектор, а затем породить
пространственный отклик вокруг его неизбежного дефекта.
\end{nogocriterion}

\begin{protocol}
\begin{enumerate}
 \item совершенное замыкание: \textbf{обратимость};
 \item индекс ориентированной ветви: \textbf{$-15$};
 \item минимальная сумма ядра и коядра: \textbf{$15$};
 \item явное насыщение границы: \textbf{получено};
 \item минимальный нормированный дефицит: \textbf{$1/7$};
 \item устойчивость при компактных деформациях: \textbf{да};
 \item уничтожение дефекта Real-спариванием: \textbf{нет};
 \item совпадение с минимумом $S_{\rm loop}$: \textbf{точное};
 \item буквальная плотность упаковки: \textbf{не используется};
 \item выбор сектора пятнадцать самим вакуумом: \textbf{не выведен};
 \item физическая энергия и радиус: \textbf{не выведены};
 \item следующий гейт: \textbf{вакуумный отклик на неизбежный дефект}.
\end{enumerate}
\end{protocol}

Машинный сертификат сохранён в
\path{s2t_v5_topological_closure_deficit_gate_results.json}.